% from ega4-2-fr.tex
\oldpage[IV-2]{100}

Nous démontrerons plus loin (13.1.1) que, sous les conditions de (5.6.5), on
a toujours $\delta(x)\geq0$, et (5.6.5.2) précise donc dans ce cas (5.5.2.1).

\begin{corollary}[5.6.6]
\label{IV.2.s5.6.6-fr}
Sous les hypothèses générales de (5.6.5), on a
\begin{equation}
\dim(X)\leq\dim(Y)+e.
\tag{5.6.6.1}\label{IV.2.s5.6.6.1-fr}
\end{equation}
Si l'on suppose en outre $Y$ universellement caténaire, alors, pour que les
deux membres de (5.6.6.1) soient égaux, il faut et il suffit que l'on ait
\begin{equation}
\dim(Y)=\sup_{y\in f(X)}\dim(\mathcal O_y).
\tag{5.6.6.2}\label{IV.2.s5.6.6.2-fr}
\end{equation}
En particulier, les deux membres de (5.6.6.1) sont égaux lorsque $Y$ est
localement de type fini sur un corps.
\end{corollary}

L'inégalité (5.6.6.1) résulte en effet de (5.6.5.1) appliqué en un point
fermé $x$ de $X$, compte tenu de (5.1.4.2) et de (5.1.11). D'autre part,
toute fibre non vide $f^{-1}(y)$ contient un point $x$ fermé dans cette
fibre (5.1.11), et si l'on suppose $Y$ universellement caténaire, on a donc
en ce point la relation (5.6.5.3). En prenant les bornes supérieures des deux
membres de (5.6.5.3) lorsque $x$ parcourt $X$, et en notant que tout point
$x$ fermé dans $X$ est \emph{a fortiori} fermé dans $f^{-1}(f(x))$, la
seconde assertion du corollaire résulte de (5.1.4.1) et (5.1.4.2).

Pour prouver la dernière assertion, notons qu'il y a un voisinage ouvert
affine $U$ du point générique de $X$ tel que $f|U$ soit de type fini, donc
$f(U)$ est une partie constructible de $Y$, dense dans $Y$ (1.8.5), et par
suite contient une partie ouverte non vide (donc partout dense) de $Y$
$(0_{\mathrm{III}},9.2.2)$. L'hypothèse que $Y$ est localement de type fini
sur un corps assure alors, d'une part que $Y$ est universellement caténaire
(5.6.4), et d'autre part que les deux membres de (5.6.6.2) sont égaux
(5.2.2 et 4.1.1.3).

On notera que la condition (5.6.6.2) est trivialement vérifiée lorsque $f$
est surjectif.

\begin{corollary}[5.6.7]
\label{IV.2.s5.6.7-fr}
Soient $Y$ un préschéma localement noethérien, $f:X\to Y$ un morphisme
localement de type fini, $n$ un entier. Si, pour tout $y\in Y$,
$\dim(f^{-1}(y))\leq n$, alors on a
\begin{equation}
\dim(X)\leq\dim(Y)+n.
\tag{5.6.7.1}\label{IV.2.s5.6.7.1-fr}
\end{equation}
\end{corollary}

En vertu de (0, 14.1.7) et de (1.5.4), on peut se borner au cas où $X$ et
$Y$ sont affines, donc $f$ de type fini, $X$ et $Y$ réduits; soient $X_i$
($1\leq i\leq m$) les sous-préschémas fermés (intègres) de $X$ ayant pour
espaces sous-jacents les composantes irréductibles de $X$; on a
$\dim(X)=\sup_i(\dim(X_i))$. Si $Z_i$ est le sous-préschéma fermé réduit de
$Y$ ayant pour espace sous-jacent $\overline{f(X_i)}$, $Z_i$ est intègre
$(0_{\mathrm I},2.1.5)$; la restriction $X_i\to Y$ de $f$ se factorise en
$X_i\xrightarrow{g_i}Z_i\xrightarrow{j_i}Y$, où $j_i$ est l'injection
canonique (I, 5.2.2), et $g_i$ est dominant et de type fini (1.5.4).

